% Part II — The Big Idea: Geometry, Not Type Theory (Slides 21–40)
\providecommand{\jugmark}[1]{\textbf{\color{jugeoBlue}#1}}

% ---------------------------------------------------------------------------
% Slide 21: Section divider
% ---------------------------------------------------------------------------
\begin{frame}
\frametitle{}
\vfill
\begin{center}
  {\Huge\color{jugeoBlue}\textbf{Part II}}\\[12pt]
  {\LARGE\color{jugeoDark}The Big Idea: Geometry, Not Type Theory}\\[20pt]
  \begin{tikzpicture}
    \draw[jugeoBlue, line width=2pt] (-3,0) -- (3,0);
  \end{tikzpicture}\\[12pt]
  {\large\color{gray}\itshape Replace proof terms with landscapes of evidence}
\end{center}
\vfill
\end{frame}

% ---------------------------------------------------------------------------
% Slide 22: Two Ways to Think About Proofs
% ---------------------------------------------------------------------------
\begin{frame}
\frametitle{Two Ways to Think About Proofs}
\begin{columns}[T]
\begin{column}{0.48\textwidth}
\begin{alertblock}{Traditional (Curry-Howard)}
  A proof is a \jugmark{PROGRAM}
  \begin{itemize}
    \item A term in a typed lambda calculus
    \item A syntax tree
    \item Grammar: right or wrong
  \end{itemize}
\end{alertblock}
\end{column}
\begin{column}{0.48\textwidth}
\begin{exampleblock}{JuGeo Approach}
  A proof is a \jugmark{MAP}
  \begin{itemize}
    \item A landscape of evidence
    \item Local measurements, globally consistent
    \item Geography: explore and chart
  \end{itemize}
\end{exampleblock}
\end{column}
\end{columns}
\vspace{10pt}
\begin{center}
\begin{tikzpicture}
  \node[fill=jugeoLight, rounded corners, inner sep=8pt, text width=10cm, align=center]
    {\itshape Think: checking grammar vs.\ charting a landscape};
\end{tikzpicture}
\end{center}
\end{frame}

% ---------------------------------------------------------------------------
% Slide 23: Analogy — Weather Forecasting
% ---------------------------------------------------------------------------
\begin{frame}
\frametitle{Analogy: Weather Forecasting}
\begin{center}
\begin{tikzpicture}[
    station/.style={circle, draw=jugeoBlue, fill=jugeoLight, minimum size=8mm,
                    font=\scriptsize\bfseries, line width=1pt},
    every edge/.style={draw=jugeoBlue!50, line width=0.8pt, dashed}
  ]
  % city outline
  \draw[jugeoDark, rounded corners=12pt, line width=1.5pt, fill=jugeoLight!30]
    (-3.5,-2) rectangle (3.5,2.2);
  \node[anchor=north west, font=\footnotesize\color{jugeoDark}] at (-3.3,2.1) {City};

  % stations
  \node[station] (s1) at (-2, 1)  {22\textdegree};
  \node[station] (s2) at ( 0, 1.5){23\textdegree};
  \node[station] (s3) at ( 2, 0.8){22\textdegree};
  \node[station] (s4) at (-1.5,-0.8){21\textdegree};
  \node[station] (s5) at ( 1,-1) {22\textdegree};
  \node[station] (s6) at ( 2.5,-1.5){23\textdegree};

  % edges
  \draw (s1)--(s2); \draw (s2)--(s3); \draw (s1)--(s4);
  \draw (s4)--(s5); \draw (s5)--(s3); \draw (s5)--(s6); \draw (s3)--(s6);
\end{tikzpicture}
\end{center}
\vspace{4pt}
\begin{itemize}
  \item Each station measures \jugmark{local} temperature
  \item Neighbors \jugmark{agree} $\Rightarrow$ glue into a global temperature map
  \item This idea is the heart of JuGeo's approach to proofs
\end{itemize}
\end{frame}

% ---------------------------------------------------------------------------
% Slide 24: Your Codebase Is a Map
% ---------------------------------------------------------------------------
\begin{frame}
\frametitle{Your Codebase Is a Map}
\begin{center}
\begin{tikzpicture}[
    mod/.style={rectangle, rounded corners, draw=jugeoBlue, fill=jugeoLight,
                minimum width=22mm, minimum height=10mm, font=\small\bfseries,
                line width=1pt},
    conn/.style={-{Stealth[length=5pt]}, jugeoBlue!70, line width=0.8pt}
  ]
  \node[mod] (app)  at (0,2.4)  {myapp};
  \node[mod] (auth) at (-3,0.6) {auth};
  \node[mod] (db)   at (0,0.6)  {database};
  \node[mod] (api)  at (3,0.6)  {api};
  \node[mod] (tok)  at (-3,-1.2){tokens};
  \node[mod] (usr)  at (0,-1.2) {users};
  \node[mod] (rtr)  at (3,-1.2) {routes};

  \draw[conn] (app)--(auth);
  \draw[conn] (app)--(db);
  \draw[conn] (app)--(api);
  \draw[conn] (auth)--(tok);
  \draw[conn] (db)--(usr);
  \draw[conn] (api)--(rtr);
  \draw[conn] (rtr)--(usr);
  \draw[conn] (tok)--(usr);
\end{tikzpicture}
\end{center}
\vspace{6pt}
Each function, module, or class is a \jugmark{location} on the map.\\[4pt]
JuGeo calls these locations \jugmark{coordinates}.
\end{frame}

% ---------------------------------------------------------------------------
% Slide 25: What Is a Coordinate?
% ---------------------------------------------------------------------------
\begin{frame}
\frametitle{What Is a Coordinate?}
\begin{block}{Definition}
  A \jugmark{coordinate} is a specific place in your codebase:\\[2pt]
  \texttt{module\,.\,function\,.\,line}
\end{block}
\vspace{8pt}
\begin{exampleblock}{Example}
  \texttt{myapp.auth.validate\_token}\\[4pt]
  is a coordinate pointing to the \texttt{validate\_token} function\\
  inside the \texttt{auth} module of \texttt{myapp}.
\end{exampleblock}
\vspace{8pt}
\begin{center}
\begin{tikzpicture}
  \node[fill=jugeoLight, rounded corners, inner sep=6pt, font=\footnotesize]
    {Like GPS coordinates --- but for code};
\end{tikzpicture}
\end{center}
\end{frame}

% ---------------------------------------------------------------------------
% Slide 26: What Are Covers?
% ---------------------------------------------------------------------------
\begin{frame}
\frametitle{What Are Covers?}
\begin{columns}[T]
\begin{column}{0.55\textwidth}
\begin{block}{Definition}
  A \jugmark{cover} is a collection of test inputs that\\
  together cover the behavior space.
\end{block}
\vspace{6pt}
\begin{exampleblock}{Example}
  10 test inputs that exercise\\
  \textbf{all code paths} of a function.
\end{exampleblock}
\end{column}
\begin{column}{0.42\textwidth}
\begin{center}
\begin{tikzpicture}[
    dot/.style={circle, fill=jugeoGreen, inner sep=2pt}
  ]
  \draw[jugeoDark, rounded corners=8pt, line width=1pt, fill=jugeoLight!30]
    (-1.4,-1.4) rectangle (1.4,1.4);
  \node[font=\scriptsize\color{jugeoDark}] at (0,1.65) {Behavior space};
  \foreach \x/\y in {-0.8/0.7, 0.1/1.0, 0.9/0.5, -1.0/-0.2, -0.3/0.1,
                      0.6/-0.1, -0.6/-0.9, 0.2/-0.7, 0.9/-1.0, -0.1/0.6} {
    \node[dot] at (\x,\y) {};
  }
\end{tikzpicture}\\[2pt]
{\scriptsize Enough stations $=$ full coverage}
\end{center}
\end{column}
\end{columns}
\end{frame}

% ---------------------------------------------------------------------------
% Slide 27: What Is Descent?
% ---------------------------------------------------------------------------
\begin{frame}
\frametitle{What Is Descent?}
\begin{center}
  \jugmark{Descent} $=$ gluing local truths into a global truth
\end{center}
\vspace{6pt}
\begin{center}
\begin{tikzpicture}[
    patch/.style={ellipse, draw=jugeoGreen, fill=jugeoGreen!10,
                  minimum width=28mm, minimum height=16mm,
                  line width=1pt, font=\small},
    arr/.style={-{Stealth[length=5pt]}, jugeoBlue, line width=1.2pt}
  ]
  \node[patch] (a) at (-2.5,0) {local proof $A$};
  \node[patch] (b) at ( 0,0)   {local proof $B$};
  \node[patch] (c) at ( 2.5,0) {local proof $C$};

  \node[rectangle, rounded corners, draw=jugeoBlue, fill=jugeoLight,
        minimum width=80mm, minimum height=10mm, line width=1.5pt,
        font=\bfseries] (g) at (0,-2) {Global proof};

  \draw[arr] (a.south) -- (g.north west);
  \draw[arr] (b.south) -- (g.north);
  \draw[arr] (c.south) -- (g.north east);
\end{tikzpicture}
\end{center}
\vspace{4pt}
\begin{itemize}
  \item Every test input confirms the spec \textcolor{jugeoGreen}{\checkmark}
  \item Test inputs cover everything \textcolor{jugeoGreen}{\checkmark}
  \item[$\Rightarrow$] Spec holds \jugmark{globally}
\end{itemize}
\end{frame}

% ---------------------------------------------------------------------------
% Slide 28: When Gluing Fails — Obstructions
% ---------------------------------------------------------------------------
\begin{frame}
\frametitle{When Gluing Fails --- Obstructions}
\begin{center}
\begin{tikzpicture}[
    station/.style={circle, draw, minimum size=10mm, font=\small\bfseries, line width=1pt}
  ]
  \node[station, draw=jugeoGreen, fill=jugeoGreen!10] (a) at (-2,0) {22\textdegree};
  \node[station, draw=jugeoRed,   fill=jugeoRed!10]   (b) at ( 2,0) {35\textdegree};
  \draw[jugeoRed, line width=2pt, decorate,
        decoration={zigzag, segment length=6, amplitude=2}]
    (a) -- (b)
    node[midway, above=6pt, font=\bfseries\color{jugeoRed}] {CONFLICT!};
\end{tikzpicture}
\end{center}
\vspace{4pt}
\begin{alertblock}{Obstruction}
  Neighboring proofs \textbf{contradict} each other at overlaps.
\end{alertblock}
\vspace{4pt}
\begin{itemize}
  \item Traditional tools: ``Unknown'' or ``timeout'' --- no explanation
  \item JuGeo: a structured \jugmark{mathematical object}\\
        that tells you \textbf{what} broke and \textbf{how} to fix it
\end{itemize}
\end{frame}

% ---------------------------------------------------------------------------
% Slide 29: The Core Insight
% ---------------------------------------------------------------------------
\begin{frame}
\frametitle{The Core Insight}
\begin{columns}[T]
\begin{column}{0.48\textwidth}
\begin{alertblock}{Traditional}
  \centering
  Proof $=$ one big term\\[6pt]
  \begin{tikzpicture}
    \node[rectangle, draw=jugeoRed, fill=jugeoRed!10, rounded corners,
          minimum width=30mm, minimum height=12mm, font=\small]
      {type-checks?};
  \end{tikzpicture}\\[6pt]
  \textbf{Yes} or \textbf{No}
\end{alertblock}
\end{column}
\begin{column}{0.48\textwidth}
\begin{exampleblock}{JuGeo}
  \centering
  Proof $=$ local evidence\\[6pt]
  \begin{tikzpicture}
    \node[ellipse, draw=jugeoGreen, fill=jugeoGreen!10,
          minimum width=14mm, minimum height=8mm, font=\scriptsize] (e1) at (-0.9,0) {$e_1$};
    \node[ellipse, draw=jugeoGreen, fill=jugeoGreen!10,
          minimum width=14mm, minimum height=8mm, font=\scriptsize] (e2) at (0.3,0.3) {$e_2$};
    \node[ellipse, draw=jugeoGreen, fill=jugeoGreen!10,
          minimum width=14mm, minimum height=8mm, font=\scriptsize] (e3) at (0.6,-0.4) {$e_3$};
  \end{tikzpicture}\\[6pt]
  Glued by \jugmark{descent}
\end{exampleblock}
\end{column}
\end{columns}
\vspace{10pt}
\begin{center}
\begin{tikzpicture}
  \node[fill=jugeoLight, rounded corners, inner sep=8pt, text width=10cm, align=center,
        font=\small]
    {Structured failures when gluing breaks $\rightarrow$
     \jugmark{actionable diagnostics}};
\end{tikzpicture}
\end{center}
\end{frame}

% ---------------------------------------------------------------------------
% Slide 30: The Judgment — JuGeo's Central Object
% ---------------------------------------------------------------------------
\begin{frame}
\frametitle{The Judgment --- JuGeo's Central Object}
\begin{center}
\begin{tikzpicture}
  \node[rectangle, rounded corners=8pt, draw=jugeoBlue, fill=jugeoLight,
        line width=2pt, inner sep=14pt, font=\large, text width=10cm, align=center]
    {A \jugmark{judgment} records:\\[8pt]
     \textbf{WHAT} is claimed\\
     \textbf{WHERE} it lives\\
     \textbf{WHAT} evidence supports it\\
     \textbf{HOW MUCH} you should trust it};
\end{tikzpicture}
\end{center}
\vspace{10pt}
\begin{itemize}
  \item Not a boolean (\texttt{True}/\texttt{False})
  \item Not a proof term (a lambda expression)
  \item An \jugmark{8-component record} --- the central data structure of JuGeo
\end{itemize}
\end{frame}

% ---------------------------------------------------------------------------
% Slide 31: The Judgment 8-Tuple — Overview
% ---------------------------------------------------------------------------
\begin{frame}
\frametitle{The Judgment 8-Tuple}
\begin{center}
  {\Large $\mathcal{J} = (c,\;\varphi,\;A,\;E,\;O,\;B,\;T,\;\Pi)$}
\end{center}
\vspace{6pt}
\begin{center}
\begin{tikzpicture}[
    comp/.style={rectangle, rounded corners, minimum width=28mm, minimum height=9mm,
                 draw=jugeoBlue, line width=0.8pt, font=\small, align=center}
  ]
  \node[comp, fill=jugeoBlue!10]   (c)  at (-4, 1.2) {$c$ --- Coordinate};
  \node[comp, fill=jugeoBlue!15]   (ph) at (-4, 0)   {$\varphi$ --- Proposition};
  \node[comp, fill=jugeoBlue!10]   (a)  at (-4,-1.2) {$A$ --- Carrier};
  \node[comp, fill=jugeoBlue!15]   (e)  at (-4,-2.4) {$E$ --- Evidence};
  \node[comp, fill=jugeoGreen!15]  (o)  at ( 1, 1.2) {$O$ --- Obligations};
  \node[comp, fill=jugeoRed!12]    (b)  at ( 1, 0)   {$B$ --- Obstructions};
  \node[comp, fill=jugeoOrange!15] (t)  at ( 1,-1.2) {$T$ --- Trust};
  \node[comp, fill=jugeoBlue!10]   (pi) at ( 1,-2.4) {$\Pi$ --- Provenance};

  \node[draw=jugeoBlue, rounded corners, line width=1.5pt,
        fit=(c)(ph)(a)(e)(o)(b)(t)(pi), inner sep=8pt] {};
\end{tikzpicture}
\end{center}
\vspace{4pt}
\begin{center}
  {\footnotesize We'll walk through each component next.}
\end{center}
\end{frame}

% ---------------------------------------------------------------------------
% Slide 32: Component 1 — Coordinate (c)
% ---------------------------------------------------------------------------
\begin{frame}
\frametitle{Component 1: Coordinate ($c$)}
\begin{block}{\jugmark{WHERE} in the codebase}
  A coordinate pinpoints the exact location being verified.
\end{block}
\vspace{6pt}
\begin{exampleblock}{Example}
  \texttt{myapp.auth.validate\_token.line\_42}
\end{exampleblock}
\vspace{6pt}
\begin{center}
\begin{tikzpicture}
  \node[circle, draw=jugeoBlue, fill=jugeoBlue!15, minimum size=12mm,
        font=\footnotesize\bfseries, line width=1.5pt] (pin) {$c$};
  \draw[-{Stealth[length=5pt]}, jugeoBlue, line width=1pt]
    (pin.east) -- ++(1.5,0)
    node[right, font=\small, text width=5cm]
      {\texttt{myapp.auth.validate\_token}\\[-1pt]\footnotesize line 42};
\end{tikzpicture}
\end{center}
\vspace{6pt}
\begin{center}
  {\small Like \jugmark{GPS coordinates} --- but for code.}
\end{center}
\end{frame}

% ---------------------------------------------------------------------------
% Slide 33: Component 2 — Proposition (φ)
% ---------------------------------------------------------------------------
\begin{frame}
\frametitle{Component 2: Proposition ($\varphi$)}
\begin{block}{\jugmark{WHAT} is claimed}
  The logical statement being verified at coordinate $c$.
\end{block}
\vspace{8pt}
\begin{exampleblock}{Example}
  \begin{center}
    ``\texttt{validate\_token} returns a non-negative integer''\\[6pt]
    $\varphi \;\equiv\; \forall\, x.\;\texttt{validate\_token}(x) \geq 0$
  \end{center}
\end{exampleblock}
\vspace{8pt}
\begin{center}
\begin{tikzpicture}
  \node[fill=jugeoLight, rounded corners, inner sep=8pt, font=\small, text width=9cm,
        align=center]
    {$c$ tells you \textbf{where} to look; $\varphi$ tells you \textbf{what} to check.};
\end{tikzpicture}
\end{center}
\end{frame}

% ---------------------------------------------------------------------------
% Slide 34: Component 3 — Carrier (A)
% ---------------------------------------------------------------------------
\begin{frame}
\frametitle{Component 3: Carrier ($A$)}
\begin{block}{\jugmark{TYPE} of the claim}
  The mathematical domain in which $\varphi$ is stated.
\end{block}
\vspace{6pt}
\begin{exampleblock}{Example}
  \begin{center}
    $A = \texttt{int} \to \texttt{int}$\\[6pt]
    The function maps integers to integers.
  \end{center}
\end{exampleblock}
\vspace{8pt}
\begin{itemize}
  \item JuGeo's version of \emph{dependent types}
  \item Determines which proof methods are valid
  \item Enables the system to compose proofs across modules
\end{itemize}
\end{frame}

% ---------------------------------------------------------------------------
% Slide 35: Component 4 — Evidence Bundle (E)
% ---------------------------------------------------------------------------
\begin{frame}
\frametitle{Component 4: Evidence Bundle ($E$)}
\begin{block}{\jugmark{THE PROOF} --- from multiple channels}
  Not one proof, but a \textbf{bundle} of evidence.
\end{block}
\vspace{4pt}
\begin{center}
\begin{tikzpicture}[
    ev/.style={rectangle, rounded corners, draw, minimum width=26mm,
               minimum height=9mm, font=\small, line width=0.8pt}
  ]
  \node[ev, draw=jugeoBlue, fill=jugeoBlue!10]   (z3)  at (-3,0) {Z3 solver proof};
  \node[ev, draw=jugeoGreen, fill=jugeoGreen!10]  (rt)  at (0,0)  {Runtime test};
  \node[ev, draw=jugeoOrange, fill=jugeoOrange!10] (ai) at (3,0)  {AI suggestion};

  \node[rectangle, rounded corners, draw=jugeoBlue, fill=jugeoLight,
        minimum width=24mm, minimum height=10mm, font=\bfseries,
        line width=1.5pt] (e) at (0,-1.8) {$E$};

  \draw[-{Stealth[length=4pt]}, jugeoBlue, line width=0.8pt] (z3.south) -- (e);
  \draw[-{Stealth[length=4pt]}, jugeoGreen, line width=0.8pt] (rt.south) -- (e);
  \draw[-{Stealth[length=4pt]}, jugeoOrange, line width=0.8pt] (ai.south) -- (e);
\end{tikzpicture}
\end{center}
\vspace{4pt}
\begin{center}
  {\small Like having \jugmark{multiple witnesses} in court---\\[2pt]
   each adds confidence from a different angle.}
\end{center}
\end{frame}

% ---------------------------------------------------------------------------
% Slide 36: Component 5 — Obligations (O)
% ---------------------------------------------------------------------------
\begin{frame}
\frametitle{Component 5: Obligations ($O$)}
\begin{block}{\jugmark{TODO list} --- what still needs proving}
  Obligations are proof goals that remain open.
\end{block}
\vspace{6pt}
\begin{center}
\begin{tikzpicture}[
    ob/.style={rectangle, rounded corners=3pt, font=\small, minimum width=60mm,
               minimum height=7mm, align=left, inner xsep=6pt}
  ]
  \node[ob, draw=jugeoGreen, fill=jugeoGreen!10]
    (o1) at (0, 0.9) {\textcolor{jugeoGreen}{\checkmark}~Arithmetic bounds --- \emph{discharged by Z3}};
  \node[ob, draw=jugeoGreen, fill=jugeoGreen!10]
    (o2) at (0, 0)   {\textcolor{jugeoGreen}{\checkmark}~Edge case $x=0$ --- \emph{discharged by test}};
  \node[ob, draw=jugeoOrange, fill=jugeoOrange!10]
    (o3) at (0,-0.9) {\textcolor{jugeoOrange}{?}~Concurrency safety --- \emph{open}};
\end{tikzpicture}
\end{center}
\vspace{6pt}
\begin{itemize}
  \item Some discharged by \jugmark{Z3}, some by \jugmark{tests}
  \item Open obligations are tracked, not forgotten
  \item JuGeo never silently drops proof goals
\end{itemize}
\end{frame}

% ---------------------------------------------------------------------------
% Slide 37: Component 6 — Obstructions (B)
% ---------------------------------------------------------------------------
\begin{frame}
\frametitle{Component 6: Obstructions ($B$)}
\begin{block}{Persistent \jugmark{blockers}}
  When something \textbf{can't} be proved, the failure itself becomes data.
\end{block}
\vspace{4pt}
\begin{columns}[T]
\begin{column}{0.48\textwidth}
\begin{alertblock}{F* / Lean}
  \centering
  \texttt{Unknown}\\[4pt]
  \texttt{(timeout after 30s)}\\[8pt]
  {\footnotesize No explanation.}
\end{alertblock}
\end{column}
\begin{column}{0.48\textwidth}
\begin{exampleblock}{JuGeo}
  \centering
  \texttt{Obstruction:}\\[2pt]
  {\footnotesize overlap at \texttt{auth$\cap$db}}\\
  {\footnotesize type mismatch \texttt{int/str}}\\[4pt]
  {\footnotesize \jugmark{What} broke $+$ \jugmark{how} to fix it.}
\end{exampleblock}
\end{column}
\end{columns}
\vspace{6pt}
\begin{center}
  {\small Obstructions are \jugmark{first-class mathematical objects},\\
   not error messages.}
\end{center}
\end{frame}

% ---------------------------------------------------------------------------
% Slide 38: Component 7 — Trust (T)
% ---------------------------------------------------------------------------
\begin{frame}
\frametitle{Component 7: Trust ($T$)}
\begin{block}{\jugmark{HOW MUCH} to trust this judgment}
  Not binary! A 7-level algebra.
\end{block}
\vspace{4pt}
\begin{center}
\begin{tikzpicture}[
    lev/.style={rectangle, rounded corners=3pt, minimum width=42mm,
                minimum height=7mm, font=\footnotesize\bfseries, align=center,
                line width=0.6pt}
  ]
  \node[lev, fill=jugeoRed!70, text=white]    at (0, 2.4) {CONTRADICTED};
  \node[lev, fill=jugeoRed!40, text=white]    at (0, 1.6) {REFUTED};
  \node[lev, fill=jugeoOrange!40]              at (0, 0.8) {UNCERTAIN};
  \node[lev, fill=jugeoOrange!25]              at (0, 0)   {PLAUSIBLE};
  \node[lev, fill=jugeoGreen!25]               at (0,-0.8) {LIKELY};
  \node[lev, fill=jugeoGreen!45]               at (0,-1.6) {VERIFIED\_TEST};
  \node[lev, fill=jugeoGreen!70, text=white]   at (0,-2.4) {VERIFIED\_PROOF};

  \draw[-{Stealth[length=6pt]}, line width=1.2pt, jugeoBlue]
    (3, 2.4) -- (3,-2.4) node[midway, right, font=\small] {increasing trust};
\end{tikzpicture}
\end{center}
\vspace{4pt}
\begin{center}
  {\footnotesize Deep dive in \textbf{Part III}.}
\end{center}
\end{frame}

% ---------------------------------------------------------------------------
% Slide 39: Component 8 — Provenance (Π)
% ---------------------------------------------------------------------------
\begin{frame}
\frametitle{Component 8: Provenance ($\Pi$)}
\begin{block}{\jugmark{WHO} produced this judgment and \jugmark{HOW}}
  Full audit trail for every piece of evidence.
\end{block}
\vspace{6pt}
\begin{center}
\begin{tikzpicture}[
    entry/.style={rectangle, rounded corners=3pt, draw=jugeoBlue,
                  fill=jugeoLight, minimum width=80mm, minimum height=8mm,
                  font=\small, align=left, inner xsep=6pt, line width=0.6pt}
  ]
  \node[entry] (p1) at (0, 0.9)
    {\texttt{14:30:01}~Z3 proved arithmetic bounds};
  \node[entry] (p2) at (0, 0)
    {\texttt{14:30:02}~pytest confirmed edge case $x{=}0$};
  \node[entry] (p3) at (0,-0.9)
    {\texttt{14:30:03}~GPT-4 suggested invariant (unverified)};

  \draw[-{Stealth[length=4pt]}, jugeoBlue!50, line width=0.6pt]
    (-4.5,0.9)--(-4.5,-0.9) node[midway, left, font=\scriptsize, rotate=90, anchor=south]
    {time};
\end{tikzpicture}
\end{center}
\vspace{6pt}
\begin{itemize}
  \item Every evidence channel is \jugmark{timestamped} and \jugmark{attributed}
  \item Enables reproducibility and blame tracking
\end{itemize}
\end{frame}

% ---------------------------------------------------------------------------
% Slide 40: Putting It All Together
% ---------------------------------------------------------------------------
\begin{frame}
\frametitle{Putting It All Together}
\begin{center}
  {\small A judgment for \texttt{validate\_token}:}
\end{center}
\vspace{2pt}
\begin{center}
\begin{tikzpicture}[
    row/.style={rectangle, rounded corners=2pt, minimum width=100mm,
                minimum height=7.5mm, font=\small, align=left, inner xsep=6pt,
                line width=0.5pt, anchor=west},
    lab/.style={font=\small\bfseries\color{jugeoBlue}, anchor=east}
  ]
  \def\yy{0.85}
  \node[lab] at (-0.2, 3*\yy) {$c$};
  \node[row, draw=jugeoBlue, fill=jugeoBlue!6]  at (0, 3*\yy)
    {\texttt{myapp.auth.validate\_token}};

  \node[lab] at (-0.2, 2*\yy) {$\varphi$};
  \node[row, draw=jugeoBlue, fill=jugeoBlue!6]  at (0, 2*\yy)
    {$\forall\, x.\;\texttt{validate\_token}(x) \geq 0$};

  \node[lab] at (-0.2, 1*\yy) {$A$};
  \node[row, draw=jugeoBlue, fill=jugeoBlue!6]  at (0, 1*\yy)
    {\texttt{str} $\to$ \texttt{int}};

  \node[lab] at (-0.2, 0*\yy) {$E$};
  \node[row, draw=jugeoGreen, fill=jugeoGreen!6] at (0, 0*\yy)
    {Z3 proof $+$ 42 passing tests $+$ GPT-4 invariant};

  \node[lab] at (-0.2,-1*\yy) {$O$};
  \node[row, draw=jugeoOrange, fill=jugeoOrange!6] at (0,-1*\yy)
    {concurrency safety (open)};

  \node[lab] at (-0.2,-2*\yy) {$B$};
  \node[row, draw=jugeoGreen, fill=jugeoGreen!6] at (0,-2*\yy)
    {$\varnothing$ --- no obstructions};

  \node[lab] at (-0.2,-3*\yy) {$T$};
  \node[row, draw=jugeoGreen, fill=jugeoGreen!6] at (0,-3*\yy)
    {\textsc{verified\_test} (level 6 of 7)};

  \node[lab] at (-0.2,-4*\yy) {$\Pi$};
  \node[row, draw=jugeoBlue, fill=jugeoBlue!6]  at (0,-4*\yy)
    {Z3 14:30, pytest 14:31, GPT-4 14:32};
\end{tikzpicture}
\end{center}
\vspace{2pt}
\begin{center}
  {\small This is what \jugmark{JuGeo produces} for every verified location.}
\end{center}
\end{frame}
